% Prevodjenje: pdflatex SystemDescription_docs.tex (dva puta za sadrzaj)
% Fajl je samostalan: nisu potrebne spoljne slike niti bibliografija.
\documentclass[a4paper,10pt]{article}
\usepackage[utf8]{inputenc}
\usepackage[T1]{fontenc}
\usepackage[serbian]{babel}
\usepackage{amssymb}
\usepackage{amsmath}
\usepackage{enumitem}
\usepackage{graphicx}
\usepackage[hidelinks]{hyperref}

\title{Korišćenje izlaznih fajlova Valgrind-ovog alata Massif u pronalasku neispravnog/nepotrebnog korišćenja memorije}
\author{Marko Nikitović\\1007/2024 \and Vladeta Vujačić\\1017/2024}
\date{Septembar 2026}

\begin{document}
\maketitle

\begin{abstract}
U ovom radu razvijena je desktop aplikacija za pregled i analizu Valgrind Massif profila.
Prikazuje potrošnju memorije kroz vreme i izdvaja rast sa zadržavanjem memorije,
memoriju prisutnu na kraju profila, nagle skokove i visok dodatni trošak
alokatora. Uz nalaze prikazuje izmerene vrednosti i, kada ima dovoljno podataka,
mesto alokacije u izvornom kodu. U radu su opisani osnovni moduli aplikacije,
pravila analize i ograničenja zaključaka koji se mogu izvesti iz profila.
\end{abstract}

\tableofcontents

\section{Uvod i opis problema}

Program može da daje tačne rezultate, a da pritom nepotrebno zadržava memoriju
ili pravi veliki broj sitnih alokacija. Sam podatak o najvećoj potrošnji memorije
ne pokazuje da li je ona rasla tokom celog izvršavanja, naglo skočila u jednom
trenutku ili ostala zauzeta do kraja.

Valgrind Massif beleži potrošnju dinamički alocirane memorije, odnosno hipa
(engl.\ heap), u nizu snimaka stanja. Snimak sadrži položaj u toku izvršavanja,
veličinu korisnog hipa i procenu dodatne memorije potrebne alokatoru. Neki snimci
sadrže i stablo poziva povezano sa alokacijama. Ručno pregledanje takvog
tekstualnog profila otežava praćenje promena i izdvajanje mesta koja treba
proveriti u izvornom kodu.

Aplikacija razvijena u okviru rada učitava ove profile, prikazuje ih grafički i izdvaja četiri obrasca:
postepen rast sa visokim zadržavanjem memorije, memoriju prisutnu u poslednjem
snimku, nagli skok potrošnje i visok dodatni trošak alokatora. Svaki nalaz sadrži
objašnjenje, izmerene vrednosti i predlog za dalju proveru. Cilj je da se
programeru olakša pregled konkretnog izvršavanja, a ne da se dokaže ispravnost
upravljanja memorijom. Massif ne pokazuje da li su objekti i dalje dostupni preko
pokazivača, pa se iz samog profila ne može potvrditi curenje memorije.

\section{Arhitektura sistema}

Aplikacija je napisana u jeziku C\# za .NET~10. Avalonia služi za desktop
interfejs, a ScottPlot za crtanje grafikona. Sistem čine tri osnovna dela:
modeli podataka, servisi za obradu i korisnički interfejs. Modeli predstavljaju
podatke koji se razmenjuju, servisi ih učitavaju i analiziraju, a interfejs
omogućava izbor ulaza i pregled rezultata. Servisi su ovde klase unutar
aplikacije, koje imaju zasebne zadatke i izvršavaju se na istom računaru.

\subsection{Modeli podataka}

Osnovni model je profil, on sadrži podatke o izvršavanju i niz snimaka stanja
memorije. Svaki snimak čuva zauzeće memorije i, ukoliko postoji, stablo alokacija.
Rezultati analize predstavljeni su nalazima i rangiranim funkcijama. Nalaz
sadrži opis uočenog obrasca, ozbiljnost, vrednosti koje ga potkrepljuju i vezu
ka relevantnom snimku. Ovi tipovi nalaze se u direktorijumu \texttt{Models}.

\subsection{Servisi za učitavanje i analizu}

Servisi u direktorijumu \texttt{Services} obavljaju sledeće zadatke:

\begin{description}
  \item[Učitavanje profila (\texttt{MassifParser}).]
  Pretvara tekstualni fajl \texttt{massif.out.*} u model profila. Čita podatke
  o izvršavanju, snimke i stabla alokacija, tako da ostali delovi aplikacije
  mogu da rade sa tim podacima bez ponovnog tumačenja teksta.

  \item[Profilisanje C programa (\texttt{SourceProfiler}).]
  Omogućava korisniku da umesto gotovog profila izabere C fajl.
  \texttt{CompilerRunner} pokreće GCC i prevodi program, a
  \texttt{MassifRunner} pokreće njegovo izvršavanje pod Massifom.
  Zajednička pomoćna klasa \texttt{ExternalTool} pokreće spoljne procese i
  prikuplja njihov izlaz. Dobijeni profil učitava se kroz \texttt{MassifParser},
  a privremeni fajlovi se zatim uklanjaju.

  \item[Prepoznavanje obrazaca (\texttt{DetectionEngine}).]
  Prima profil i zadate pragove, priprema podatke za analizu i pokreće četiri
  pravila: LEAK, ATEXIT, SPIKE i FRAG. Vraća listu nalaza poređanih po
  ozbiljnosti. Svako pravilo je izdvojeno u posebnu klasu, dok su zajednički
  proračuni i pragovi objedinjeni. Time se izbegava ponavljanje iste obrade
  u više pravila.

  \item[Rangiranje funkcija (\texttt{HotFunctionAnalyzer}).]
  Grupisanjem mesta alokacije izdvaja funkcije i rangira ih prema najvećem
  zabeleženom zauzeću. Ovaj pregled pomaže korisniku da pronađe velike
  potrošače memorije čak i kada nijedno pravilo ne prijavi nalaz.
\end{description}

\subsection{Korisnički interfejs i povezivanje modula}

Ulazna tačka programa pokreće Avalonia okruženje, a klasa \texttt{App} kreira
glavni prozor \texttt{MainWindow}. Glavni prozor povezuje rad servisa:
na zahtev korisnika poziva parser ili servis za profilisanje, dobijeni profil
prosleđuje servisima za prepoznavanje obrazaca i rangiranje funkcija, a zatim
prikazuje rezultate.

Tok podataka je zato isti nakon oba načina unosa, odnosno učitani profil postaje ulaz
za analizu i grafički prikaz. Prozor prikazuje graf potrošnje, snimke,
stabla alokacija, rangirane funkcije i nalaze. Izbor nalaza vodi na snimak koji
ga potkrepljuje. Posebni prozori omogućavaju promenu podešavanja i pregled
izvornog koda kada je dostupan.

Ovakva podela odvaja proračune od prikaza, tj. parser i analiza ne zavise od
Avalonia interfejsa, dok glavni prozor obrađuje korisničke radnje i osvežava
prikaz. Oba načina unosa koriste istu analizu, pa nije potrebna posebna obrada
za upravo generisan profil. Otvaranje postojećih profila ne zahteva GCC ni
Valgrind, oni su potrebni samo za direktno profilisanje C programa.

\section{Postupak analize i pravila}

Za svaki snimak sa indeksom $i$ izdvajaju se korisni hip $H_i$, dodatni trošak
alokatora $E_i$ i položaj $t_i$ u izvršavanju. Računa se maksimum
$P=\max_i H_i$, a položaji se normalizuju na interval $[0,1]$ prema prvom i
poslednjem snimku. Ako svi imaju istu vremensku oznaku, koriste se njihovi
relativni indeksi. Normalizacija omogućava da se prozori posmatranja izraze kao
deo trajanja profila, nezavisno od izabrane jedinice. Podrazumevana jedinica pri
profilisanju jeste broj izvršenih instrukcija, a ne broj sekundi.

Pravila se primenjuju nezavisno. Sa podrazumevanim podešavanjima nijedno pravilo ne daje
nalaz ako profil ima manje od deset snimaka ili ako je maksimalni hip jednak nuli.
U nastavku su navedeni podrazumevani pragovi.

\subsection{LEAK - rast sa zadržavanjem memorije}
Računa se Spirmanova korelacija između položaja u izvršavanju i zauzeća hipa.
Ona poredi rangove vrednosti, pa ne zahteva da rast bude linearan, takodje jednake
vrednosti dobijaju prosečan rang. Nalaz se daje kada istovremeno važi:
\[
  \rho(t,H)>0{,}90, \qquad
  \frac{H_{n-1}-H_0}{P}>0{,}20, \qquad
  \frac{H_{n-1}}{P}>0{,}80.
\]
Tako se traži izražen rast, dovoljno velika ukupna promena i visok nivo
zadržavanja na kraju. Sam uzlazni trend ne bi bio dovoljan jer bi mogao da predstavlja
malo povećanje ili memoriju koja je kasnije oslobođena. Nalaz je upozorenje na
moguće curenje, uz napomenu da zadržavanje može biti namerno.

Preporučuje se provera da li se memorija oslobađa kada više nije potrebna,
uz dodatnu proveru curenja pomoću Memcheck-a.

\subsection{ATEXIT - memorija u poslednjem snimku}
Ako poslednji snimak sadrži više od $10\%$ maksimalnog hipa, daje se informativni
nalaz. Ovo pravilo ne zahteva prethodni rast. Poslednji snimak se posmatra kao
poslednji dostupan podatak, a ne kao dokaz da je memorija ostala posle svih
završnih operacija programa.

Preporučuje se provera da li memorija mora da ostane zauzeta do kraja programa,
kao i dodatna provera curenja pomoću Memcheck-a.

\subsection{SPIKE - nagli skok}
Za dva susedna snimka traži se porast koji prelazi oba praga:
\[
  H_i-H_{i-1}>\max(0{,}50H_{i-1},\;0{,}10P).
\]
Relativni prag izdvaja naglu promenu, a prag u odnosu na maksimum sprečava da se
mali porast sa skoro praznog hipa proglasi značajnim. Uzastopni skokovi spajaju
se u jedan događaj. Zatim se traži najmanje zauzeće u kasnijim snimcima unutar
prozora dužine $20\%$ celog profila, merenog od kraja skoka.
Ako hip padne na početni nivo uvećan za najviše $25\%$ ukupnog skoka, događaj
se označava kao prolazan i daje se informativni nalaz. U suprotnom se daje
upozorenje o zadržanom povećanju u tom prozoru. Ako nema kasnijih snimaka unutar
prozora, ishod ostaje nepoznat i nalaz je informativan.

Ako prolazni skok dovodi do prevelike potrošnje, preporučuje se obrada
podataka u manjim delovima. Kod zadržanog povećanja treba proveriti da li je
dodatna memorija i dalje potrebna i razmotriti ograničavanje njenog zauzeća.
Za nepoznat ishod preporučuje se prikupljanje dodatnih snimaka nakon skoka.

\subsection{FRAG - visok trošak alokatora}
Za snimke sa nenultim hipom računa se odnos $E_i/H_i$. Upozorenje se daje kada
je medijana tih odnosa veća od $0{,}30$. Medijana smanjuje uticaj pojedinačnih
ekstremnih odnosa, naročito pri malom zauzeću hipa. Ovaj nalaz može da ukaže na
veliki broj sitnih alokacija. Uprkos nazivu pravila, on nije direktno merenje
fragmentacije (koristi Massif-ovu procenu troška po bloku i poravnanja).

Preporučuje se provera velikog broja sitnih alokacija i, kada je moguće,
njihovo grupisanje u veće blokove.

\subsection{Povezivanje nalaza sa izvornim kodom}

Mesta alokacije uzimaju se iz neposrednih potomaka korena stabla. Ponovljene
oznake se sabiraju, a zbirne stavke oblika \texttt{in N places} preskaču, jer ne
predstavljaju jedno poznato mesto. Zato njihov zbir ne mora biti jednak ukupnom
hipu.

Za LEAK se mesto navodi samo ako prvi dostupan snimak sa stablom pripada prvoj,
a poslednji poslednjoj četvrtini profila. Bira se najveći zabeleženi porast koji
objašnjava najmanje $10\%$ ukupnog rasta. ATEXIT koristi isključivo stablo
poslednjeg snimka i izdvaja najveće mesto koje zauzima više od $10\%$ završnog
hipa. SPIKE zahteva stabla neposredno pre i na kraju događaja, uz jedno mesto
koje objašnjava najmanje polovinu porasta. FRAG ne pripisuje nalaz konkretnoj
funkciji jer je dodatni trošak dostupan samo za ceo hip. Kada podaci nisu
dovoljni, nalaz ostaje bez imenovanog mesta alokacije.

Kada je profil generisan iz C fajla kroz aplikaciju i sumnjiva lokacija sadrži
prepoznatljiv broj reda u tom fajlu, dugme \textit{Show in source} otvara izvorni
kod i označava odgovarajući red. Pri otvaranju postojećeg profila lokacija se
može prikazati u nalazu, ali direktan pregled izvornog koda nije omogućen.

Posle primene pravila nalazi se sortiraju po ozbiljnosti. Ako LEAK i ATEXIT
izdvoje isto mesto, ATEXIT dobija dodatno obrazloženje, bez povećanja ozbiljnosti.
Rangiranje funkcija prikazuje njihove pojedinačne maksimume, koji mogu poticati
iz različitih snimaka, pa njihovi udeli zato nisu raspodela jednog zajedničkog
stanja memorije.

Pragovi su heuristički i mogu se menjati u prozoru za podešavanja. Izmene važe
tokom tekuće sesije i koriste se pri sledećem učitavanju ili generisanju profila.
Već prikazani nalazi se ne preračunavaju.

\section{Zaključak}

U ovom radu povezan je grafički pregled profila sa pravilima koja izdvajaju
sumnjive promene u potrošnji memorije. Korisnik može da pređe sa nalaza na
relevantni snimak i proveri podatke na kojima se nalaz zasniva.
Rezultat zavisi i od dostupnih snimaka, pa
kratka promena između njih može ostati nezabeležena. Zbog toga uz nalaze ostaju
konkretne vrednosti i ograničenja, kako bi korisnik mogao da proceni da li obrazac
odgovara očekivanom ponašanju programa.

\end{document}
